\lecture{6}{2026-09-29}{Miller--Rabin Primality Test}

\section*{The Miller--Rabin Primality Test}

The Miller--Rabin test strengthens the Fermat test by exploiting the fact that over a field, the equation \(x^2 = 1\) has only the trivial roots \(x = \pm 1\).

\begin{lemma}\label{lem:roots-of-unity}
    If \(p\) is prime and \(x \in \mathbb{Z}_p\) satisfies \(x^2 \equiv 1 \pmod{p}\), then \(x \equiv \pm 1 \pmod{p}\).
\end{lemma}

\begin{proof}
    The condition \(x^2 \equiv 1 \pmod{p}\) is equivalent to \(p \mid (x^2 - 1) = (x - 1)(x + 1)\).
    By Euclid's lemma, since \(p\) is prime, \(p \mid (x - 1)\) or \(p \mid (x + 1)\), which means \(x \equiv 1 \pmod{p}\) or \(x \equiv -1 \pmod{p}\).
\end{proof}

If \(n\) is composite and has at least two distinct odd prime factors, the Chinese Remainder Theorem implies that \(x^2 \equiv 1 \pmod{n}\) has at least \(2^2 = 4\) distinct solutions modulo \(n\), creating \emph{non-trivial square roots of \(1\)}.
The Miller--Rabin test hunts for such non-trivial roots in the sequence of powers leading up to \(a^{n-1}\).

Let \(n \geq 3\) be an odd integer. We factor \(n - 1\) as
\[
    n - 1 = 2^e k, \quad \text{where \(k\) is odd and \(e \geq 1\)}.
\]
For any \(a \in \mathbb{Z}_n^*\), consider the sequence obtained by repeated squaring:
\[
    a^k, \; a^{2k}, \; a^{2^2 k}, \; \dots, \; a^{2^{e-1} k}, \; a^{2^e k} = a^{n-1} \pmod{n}.
\]

\begin{theorem}[Miller's Condition~\cite{Miller1976}]\label{thm:miller-condition}
    Let \(n \geq 3\) be an odd prime and write \(n - 1 = 2^e k\) with \(k\) odd.
    For every \(a \in \mathbb{Z}_n^*\), at least one of the following holds:
    \begin{enumerate}[label=\textup{(\roman*)}]
        \item \(a^k \equiv 1 \pmod{n}\);
        \item there exists some \(r \in \{0, 1, \dots, e-1\}\) such that \(a^{2^r k} \equiv -1 \pmod{n}\).
    \end{enumerate}
\end{theorem}

\begin{proof}
    Since \(n\) is prime, Fermat's Little Theorem guarantees that the last term of the sequence is \(a^{2^e k} = a^{n-1} \equiv 1 \pmod{n}\).
    If \(a^k \equiv 1 \pmod{n}\), condition (i) holds.
    Otherwise, let \(r \in \{0, 1, \dots, e-1\}\) be the largest index such that \(a^{2^r k} \not\equiv 1 \pmod{n}\).
    Then the next term in the sequence is \(a^{2^{r+1} k} = (a^{2^r k})^2 \equiv 1 \pmod{n}\).
    Thus, \(x = a^{2^r k}\) is a square root of \(1\) modulo \(n\) that is not congruent to \(1\).
    By \Cref{lem:roots-of-unity}, the only square roots of \(1\) modulo a prime are \(\pm 1\), so we must have \(a^{2^r k} \equiv -1 \pmod{n}\), satisfying condition (ii).
\end{proof}

\begin{definition}[Strong Witness and Strong Liar]
    Let \(n \geq 3\) be an odd composite integer with \(n - 1 = 2^e k\) (\(k\) odd), and let \(a \in \mathbb{Z}_n^*\).
    \begin{itemize}
        \item We say \(a\) is a \emph{strong witness} for \(n\) if
        \[
            a^k \not\equiv 1 \pmod{n} \quad \text{and} \quad a^{2^r k} \not\equiv -1 \pmod{n} \quad \text{for all } r \in \{0, 1, \dots, e - 1\}.
        \]
        \item Otherwise, \(a\) is called a \emph{strong liar} for \(n\) (and \(n\) is said to be a \emph{strong pseudoprime} to base \(a\)).
    \end{itemize}
\end{definition}

If \(a\) is a strong witness, then either \(a^{n-1} \not\equiv 1 \pmod{n}\) (a Fermat witness) or we encounter some \(x \not\equiv \pm 1 \pmod{n}\) with \(x^2 \equiv 1 \pmod{n}\), which yields a non-trivial square root of \(1\) modulo \(n\).
In either case, \(n\) is definitively composite.

\begin{algorithm}[htbp]
    \caption{Miller--Rabin Primality Test}\label{alg:miller-rabin}
    \begin{algorithmic}[1]
        \Input odd integer \(n \geq 3\) and candidate base \(a \in \{2, \dots, n-2\}\)
        \Output \textup{composite} or \textup{probably prime}
        \State factor \(n - 1 = 2^e k\) with \(k\) odd
        \State set \(x \gets a^k \pmod{n}\)
        \If{\(x \equiv 1 \pmod{n}\) \textbf{or} \(x \equiv -1 \pmod{n}\)}
            \State \Return \textup{probably prime} \Comment{\(a\) is a strong liar}
        \EndIf
        \For{\(r = 1\) \textbf{to} \(e - 1\)}
            \State set \(x \gets x^2 \pmod{n}\)
            \If{\(x \equiv -1 \pmod{n}\)}
                \State \Return \textup{probably prime} \Comment{\(a\) is a strong liar}
            \EndIf
            \If{\(x \equiv 1 \pmod{n}\)}
                \State \Return \textup{composite} \Comment{found non-trivial square root of \(1\)}
            \EndIf
        \EndFor
        \State \Return \textup{composite} \Comment{\(a^{n-1} \not\equiv 1 \pmod{n}\)}
    \end{algorithmic}
\end{algorithm}

Unlike the Fermat test, there are no composite numbers that defeat the Miller--Rabin test for all bases: even Carmichael numbers possess many strong witnesses.

\begin{theorem}\label{thm:strong-liars-bound}
    Let \(n \geq 3\) be an odd composite integer.
    Then at most \(\frac{1}{2}\) of the elements in \(\mathbb{Z}_n^*\) are strong liars for \(n\).
\end{theorem}

\begin{proof}
    We divide the proof into two cases depending on whether \(n\) is a prime power.

    \paragraph{Case 1: \(n = p^m\) for some prime \(p\) and \(m \geq 2\).}
    If \(a \in \mathbb{Z}_n^*\) is a strong liar, then \(a^{n-1} \equiv 1 \pmod{n}\).
    By Euler's theorem, we also have \(a^{\phi(n)} \equiv 1 \pmod{n}\), where \(\phi(n) = p^{m-1}(p-1)\).
    The order of \(a\) in \(\mathbb{Z}_n^*\) must divide both \(n - 1 = p^m - 1\) and \(\phi(n) = p^{m-1}(p-1)\).
    Since \(\gcd(p^m - 1, p^{m-1}) = 1\),
    \[
        \gcd(p^m - 1, p^{m-1}(p-1)) = \gcd(p^m - 1, p - 1) = p - 1.
    \]
    Thus, every strong liar must satisfy \(a^{p-1} \equiv 1 \pmod{n}\).
    The set \(H = \{a \in \mathbb{Z}_n^* : a^{p-1} \equiv 1 \pmod{n}\}\) forms a subgroup of \(\mathbb{Z}_n^*\).
    To show that \(H\) is proper, consider \(a = 1 + p^{m-1}\).
    By the binomial theorem,
    \[
        (1 + p^{m-1})^{p-1} = 1 + (p - 1)p^{m-1} + \sum_{j=2}^{p-1} \binom{p-1}{j} p^{j(m-1)} \equiv 1 - p^{m-1} \not\equiv 1 \pmod{p^m},
    \]
    since \(m \geq 2\) and \(2(m-1) \geq m\).
    Thus \(1 + p^{m-1} \notin H\), so \(H\) is a proper subgroup of \(\mathbb{Z}_n^*\).
    By Lagrange's theorem, \(|H| \leq \frac{1}{2} |\mathbb{Z}_n^*|\), so at most half of the elements in \(\mathbb{Z}_n^*\) are strong liars.

    \paragraph{Case 2: \(n\) is not a prime power.}
    Since \(n\) is not a prime power and is odd, we can write \(n = p^m q\) where \(p\) is an odd prime, \(m \geq 1\), and \(\gcd(p^m, q) = 1\) with \(q > 1\).
    Let \(i_0\) be the maximal index in \(\{0, 1, \dots, e-1\}\) for which there exists an element \(a_0 \in \mathbb{Z}_n^*\) satisfying
    \[
        a_0^{2^{i_0} k} \equiv -1 \pmod{n}.
    \]
    Such an index exists because \((-1)^k = -1\) (as \(k\) is odd), so \(i = 0\) always satisfies this with \(a = -1\).
    Now define the set
    \[
        G = \{a \in \mathbb{Z}_n^* : a^{2^{i_0} k} \equiv \pm 1 \pmod{n}\}.
    \]
    The set \(G\) is a subgroup of \(\mathbb{Z}_n^*\).
    Moreover, every strong liar \(a\) lies in \(G\):
    \begin{itemize}
        \item If \(a^k \equiv 1 \pmod{n}\), then \(a^{2^{i_0} k} \equiv 1 \pmod{n}\), so \(a \in G\).
        \item If \(a^{2^r k} \equiv -1 \pmod{n}\) for some \(r \leq i_0\), then \(a^{2^{i_0} k} \equiv (a^{2^r k})^{2^{i_0 - r}} \equiv (-1)^{2^{i_0 - r}} \equiv \pm 1 \pmod{n}\), so \(a \in G\).
    \end{itemize}
    We now show that \(G\) is a proper subgroup of \(\mathbb{Z}_n^*\).
    By the Chinese Remainder Theorem, choose \(b \in \mathbb{Z}_n^*\) such that
    \[
        b \equiv a_0 \pmod{p^m} \quad \text{and} \quad b \equiv 1 \pmod{q}.
    \]
    Raising to the power \(2^{i_0} k\), we obtain
    \[
        b^{2^{i_0} k} \equiv a_0^{2^{i_0} k} \equiv -1 \pmod{p^m}, \quad \text{and} \quad b^{2^{i_0} k} \equiv 1^{2^{i_0} k} \equiv 1 \pmod{q}.
    \]
    Since \(p^m \geq 3\) and \(q \geq 3\), \(-1 \not\equiv 1\) modulo \(p^m\) and modulo \(q\).
    For \(b^{2^{i_0} k}\) to equal \(1 \pmod{n}\), it would have to be \(1 \pmod{p^m}\); for it to equal \(-1 \pmod{n}\), it would have to be \(-1 \pmod{q}\).
    Since neither holds, \(b^{2^{i_0} k} \not\equiv \pm 1 \pmod{n}\), and hence \(b \notin G\).
    Therefore \(G\) is a proper subgroup of \(\mathbb{Z}_n^*\), and by Lagrange's theorem,
    \[
        |G| \leq \frac{1}{2} |\mathbb{Z}_n^*|.
    \]
    Since all strong liars belong to \(G\), at most \(\frac{1}{2}|\mathbb{Z}_n^*|\) elements are strong liars.
\end{proof}

\begin{remark}[Rabin's Bound]
    A more refined analysis due to Rabin~\cite{Rabin1980} shows that for any odd composite \(n\), the number of strong liars is actually bounded by
    \[
        |\{\text{strong liars in } \mathbb{Z}_n^*\}| \leq \frac{1}{4} \phi(n).
    \]
\end{remark}

\begin{remark}[Error Probability and Complexity]
    If \(n\) is composite and we run the Miller--Rabin test with \(t\) independent, uniformly chosen bases \(a_1, \dots, a_t \in \{2, \dots, n-2\}\), the probability that \(n\) is declared ``probably prime'' in all \(t\) iterations is at most
    \[
        \left( \frac{1}{4} \right)^t = 2^{-2t}.
    \]
    Setting \(t = 64\) guarantees a false positive probability less than \(2^{-128}\), which is cryptographically negligible.
    
    Each round requires computing \(a^k \pmod{n}\) and at most \(e - 1\) modular squarings, for a total of \(O(\log n)\) modular multiplications.
    With standard \(O(\log^2 n)\) arithmetic, each round runs in \(O(\log^3 n)\) time.
    Testing \(t\) independent bases thus has bit complexity \(O(t \log^3 n)\).
\end{remark}

\section*{Languages}

A \emph{language} is a set of binary strings.

\begin{example}
    \[
        \mathsf{PRIMES} = \{x \in \{0, 1\}^* : x \text{ is the binary representation of a prime number}\}
    \]
\end{example}

\begin{definition}
    A language \(L\) is in \(\mathsf{RP}\) (randomized polynomial time) if there exists a probabilistic polynomial-time algorithm \(A\) such that for every input \(x\):
    \begin{itemize}
        \item if \(x \in L\), then \(\Pr[A(x) = \text{accept}] \geq 1-\delta\), where \(\delta < 1 - \frac{1}{\text{poly}(|x|)}\); and
        \item if \(x \notin L\), then \(\Pr[A(x) = \text{accept}] = 0\).
    \end{itemize}
\end{definition}

The Miller--Rabin test shows that the complement \(\overline{\mathsf{PRIMES}} \in \mathsf{RP}\), since if \(n\) is composite, the algorithm returns ``composite'' with probability at least \(1 - 2^{-2t}\) after \(t\) independent iterations, and if \(n\) is prime, the algorithm never returns ``composite''.

\begin{definition}
    A language \(L\) is in \(\mathsf{P}\) (deterministic polynomial time) if there exists a deterministic polynomial-time algorithm \(A\) such that for every input \(x\):
    \begin{itemize}
        \item if \(x \in L\), then \(A(x) = \text{accept}\); and
        \item if \(x \notin L\), then \(A(x) = \text{reject}\).
    \end{itemize}
\end{definition}

\begin{definition}
    A language \(L\) is in \(\mathsf{NP}\) (nondeterministic polynomial time) if there exists a polynomial-time algorithm \(A\) such that, for all \(x \in L\) there exists a polynomial-size witness \(w\) such that \(A(x, w) = \text{accept}\), and for all \(x \notin L\), there is no witness \(w\) such that \(A(x, w) = \text{accept}\).
\end{definition}

It's easy to see that \(\overline{\mathsf{PRIMES}} \in \mathsf{NP}\) as well, since if \(n\) is composite, a non-trivial factor of \(n\) serves as a polynomial-size witness that can be verified in polynomial time using the Extended Euclidean Algorithm.

A harder example is to show that \(\mathsf{PRIMES} \in \mathsf{NP}\).